\documentclass[11pt]{article}

\usepackage[margin=1in]{geometry}
\usepackage{amsmath,amssymb}
\usepackage{booktabs}
\usepackage{enumitem}
\usepackage[hidelinks]{hyperref}

\setlist[itemize]{leftmargin=1.4em}
\setlist[enumerate]{leftmargin=1.6em}

\title{SYSC 5108 Exam Study\\Assignment 3, Question 2}
\author{Jesse Levine}
\date{}

\begin{document}
\maketitle

\section*{Question}

The artificial neuron takes three inputs. For the input vector
\[
    \mathbf{d} = [0.2,\;0.5,\;0.7],
\]
calculate:

\begin{enumerate}[label=(\alph*)]
    \item the weighted sum for the neuron,
    \item the output if the activation function \(\varphi\) is a threshold
    activation with threshold \(\theta=1\),
    \item the output if \(\varphi\) is the logistic function,
    \item the output if \(\varphi\) is the rectified linear unit (ReLU) function.
\end{enumerate}

\section*{Course and Textbook Cross-References}

\begin{itemize}
    \item \textbf{Chapter 4 slides, General Neural Network:} a neuron computes
    an affine or weighted sum first, then applies a nonlinear activation function.
    \item \textbf{Chapter 4 slides, Hidden Units:} ReLU is defined as
    \[
        g(z)=\max(0,z).
    \]
    \item \textbf{Goodfellow, Bengio, and Courville, \emph{Deep Learning},
    Chapter 3:} the logistic sigmoid is
    \[
        \sigma(x)=\frac{1}{1+\exp(-x)}.
    \]
    \item \textbf{Goodfellow et al., Chapter 6:} neural-network hidden units
    usually compute an affine transformation and then apply an activation function;
    ReLU is defined as \(g(z)=\max(0,z)\).
\end{itemize}

\section*{Given Values From the Diagram}

The bias input is
\[
    d_0 = 1.
\]
The weights are:
\[
    w_0 = 0.20,\qquad
    w_1 = -0.10,\qquad
    w_2 = 0.15,\qquad
    w_3 = 0.05.
\]
The input vector gives
\[
    d_1=0.2,\qquad d_2=0.5,\qquad d_3=0.7.
\]

\section*{Step-by-Step Solution}

\subsection*{Part (a): Weighted sum}

The weighted sum is
\[
    z = w_0d_0+w_1d_1+w_2d_2+w_3d_3.
\]
Substitute the values:
\[
\begin{aligned}
    z
    &=
    (0.20)(1)+(-0.10)(0.2)+(0.15)(0.5)+(0.05)(0.7) \\
    &=
    0.20-0.020+0.075+0.035 \\
    &=
    0.290.
\end{aligned}
\]
Therefore,
\[
    \boxed{z=0.29}.
\]

\subsection*{Part (b): Threshold activation with \(\theta=1\)}

Using a threshold activation,
\[
    \varphi(z)
    =
    \begin{cases}
    1, & z\ge \theta, \\
    0, & z<\theta.
    \end{cases}
\]
Here,
\[
    z=0.29 < 1.
\]
Therefore,
\[
    \boxed{\varphi(z)=0}.
\]

\subsection*{Part (c): Logistic activation}

The logistic activation is
\[
    \varphi(z)=\sigma(z)=\frac{1}{1+\exp(-z)}.
\]
Using \(z=0.29\):
\[
\begin{aligned}
    \sigma(0.29)
    &=
    \frac{1}{1+\exp(-0.29)} \\
    &\approx
    0.5720.
\end{aligned}
\]
Therefore,
\[
    \boxed{\varphi(z)\approx 0.5720}.
\]

\subsection*{Part (d): ReLU activation}

The ReLU activation is
\[
    \varphi(z)=\max(0,z).
\]
Since \(z=0.29>0\),
\[
    \varphi(0.29)=\max(0,0.29)=0.29.
\]
Therefore,
\[
    \boxed{\varphi(z)=0.29}.
\]

\section*{Final Answer}

\[
\boxed{
\begin{array}{c|c}
\text{Quantity} & \text{Value} \\
\hline
\text{Weighted sum } z & 0.29 \\
\text{Threshold output, } \theta=1 & 0 \\
\text{Logistic output} & 0.5720 \\
\text{ReLU output} & 0.29
\end{array}
}
\]

\section*{Exam Notes}

\begin{itemize}
    \item Do not forget the bias term \(w_0d_0=(0.20)(1)\).
    \item The weighted sum \(z\) is computed once. Each activation function then
    transforms the same \(z=0.29\) in a different way.
    \item With threshold activation, compare \(z\) directly against the threshold
    \(\theta\). Here \(0.29<1\), so the output is \(0\).
    \item With ReLU, positive values pass through unchanged and negative values
    are clipped to \(0\).
\end{itemize}

\end{document}
